% !TEX TS-program = pdflatex
% !TeX program = pdflatex
% !TEX encoding = UTF-8
% !TEX spellcheck = en_US

\documentclass[xcolor=table]{beamer}


\def\accyear{2024-2025}

\institute{
	Laboratoire de la Communication dans les Systèmes Informatiques (LCSI)\\
	École  nationale Supérieure d'Informatique (ESI, ex. INI), Algiers, Algeria
}

\author[Abdelkrime Aries (\accyear)] %
{Abdelkrime Aries}

\titlegraphic{%
	\includegraphics[height=1cm]{../extra/logo/esi-logo.png}%
	\hspace*{2cm}%
	\includegraphics[height=1cm]{../extra/logo/lcsi-logo.png}%
	\hspace*{2cm}%
	\includegraphics[height=1cm]{../extra/logo/esi.nlp.pdf}%
}%\hspace*{4.75cm}~

\date{Academic year: \accyear} %\today

\title[NLP]{Natural Language Processing} 

\usefolder{../extra/beamer}

\usetheme{Karimnlp} % Antibes Boadilla Warsaw

\beamertemplatenavigationsymbolsempty



\title[ESI - NLP: 05- PoS tagging]%
{Natural Language Processing\\Chapter 05\\Part-of-speech tagging} 

\changegraphpath{../img/pos/}

\begin{document}
	
\begin{frame}
\frametitle{Natural Language Processing}
\framesubtitle{Part-of-speech tagging: Introduction}

\begin{exampleblock}{Example of two sentences}
	\begin{center}
		\Huge\bfseries
		We can can the can
		
		Will Will will the will to Will?
	\end{center}
\end{exampleblock}

\begin{itemize}
	\item Did you understand these two sentences?
	\item Find nouns, verbs, determiners, etc.
	\item \optword{can}: (1) V: have the possibility (2) V: put in a container (3) N: container
	\item \optword{will}: (1) V: express the future (2) N: masculine proper noun (3) V: determine by choice (4) N: testament
\end{itemize}

\end{frame}

\begin{frame}
\frametitle{Natural Language Processing}
\framesubtitle{Part-of-speech: Some humor}

\begin{center}
	\vgraphpage{humor/humor-pos.jpg}
\end{center}

\end{frame}

\begin{frame}
\frametitle{Natural Language Processing}
\framesubtitle{Part-of-speech: Plan}

\begin{multicols}{2}
%	\small
\tableofcontents
\end{multicols}
\end{frame}

%===================================================================================
\section{Sequences labeling}
%===================================================================================

\subsection{Description}

\begin{frame}
\frametitle{PoS tagging: Sequences labeling}
\framesubtitle{Description}


\begin{itemize}
	\item \textbf{Goal}: assigning labels to elements of a sequence.
	\item \textbf{Example}: \expword{Assigning PoS tags to a sentence's words}.
\end{itemize}

\begin{block}{Problem formulation}
	\begin{itemize}
		\item $w = w_1, \ldots, w_n$: input words sequence
		\item $t = t_1, \ldots, t_n$: output labels sequence
	\end{itemize}
	\begin{center}
		$ \arg\max\limits_t P(t | w)$
	\end{center}
	
	\begin{itemize}
		\item \optword{Generative models} $ \arg\max\limits_t P(t | w) = \arg\max\limits_t P(t) P(w | t) $
		\item \optword{Discriminative models} estimating $\arg\max\limits_t P(t | w)$ directly
	\end{itemize}
\end{block}

\end{frame}

\subsection{Applications}

\begin{frame}
\frametitle{PoS tagging: Sequences labeling}
\framesubtitle{Applications}

%eisenstein
\begin{itemize}
	\item Named Entity Recognition (NER)
	\begin{itemize}
		\item Find people, organizations, places, numbers, etc. in a sentence
	\end{itemize}
	\item PoS tagging
	\begin{itemize}
		\item Find words PoS tags (noun, verb, etc.)
	\end{itemize}
	\item Chunking
	\begin{itemize}
		\item Find different phrases in a sentence (such as nominal phrase)
	\end{itemize}
\end{itemize}


\end{frame}

\begin{frame}
\frametitle{PoS tagging: Sequences labeling}
\framesubtitle{Applications (example)}

\begin{figure}
	\centering
	\hgraphpage{exp-ner2_.pdf}
	\caption{Example of NER [\url{https://corenlp.run/}]}
\end{figure}

\begin{figure}
	\centering
	\hgraphpage{exp-pos2_.pdf}
	\caption{Example of PoS tagging [\url{https://corenlp.run/}]}
\end{figure}

\end{frame}

%\subsection{Évaluation}
%
%\begin{frame}
%\frametitle{Étiquetage morpho-syntaxique : Séquences}
%\framesubtitle{Évaluation}
%
%\begin{itemize}
%	\item 
%\end{itemize}
%
%\end{frame}

%===================================================================================
\section{Task description}
%===================================================================================

\begin{frame}
\frametitle{Part-of-speech tagging}
\framesubtitle{Task description}

\begin{itemize}
	\item \textbf{Objective}: find the PoS of each word in a sentence
	\item Define PoS
	\begin{itemize}
		\item General: \expword{N (noun), V (verb), J (adjective), A (adverb), P (preposition), Determiner (D)}
		\item Specific: \expword{NN (noun), NNP (proper noun), V (present verb), VP (past verb) ...}
	\end{itemize}
	\item Annotate a corpus using these tags
	\begin{itemize}
		\item Manually annotate a test corpus
		\item Another one for training if we want to use machine learning
	\end{itemize}
\end{itemize}

\end{frame}

\subsection{Universal classes}

\begin{frame}
\frametitle{PoS tagging: Task description}
\framesubtitle{Universal classes}

\url{https://universaldependencies.org/u/pos/}

\begin{tblr}{
		colspec = {p{.3\textwidth}p{.32\textwidth}p{.28\textwidth}},
		row{odd} = {lightblue},
		row{even} = {lightyellow},
		row{1} = {darkblue},
		rowsep = 1pt,
	} 
	\textcolor{white}{Open class} & \textcolor{white}{Close class} & \textcolor{white}{Other} \\
	
	\keyword{ADJ}:  adjective & \keyword{ADP}: adposition & \keyword{PUNCT}: punctuation \\
	\keyword{ADV}:  adverb & \keyword{AUX}: auxiliary & \keyword{SYM}: symbol \\
	\keyword{INTJ}: interjection & \keyword{CCONJ}: coordination conjunction & \keyword{X}: Other \\
	\keyword{NOUN}: noun & \keyword{DET}: determiner &  \\
	\keyword{PROPN}: proper noun & \keyword{NUM}: numerical &  \\
	\keyword{VERB}: verb & \keyword{PART}: particle &  \\
	 & \keyword{PRON}: pronoun &  \\
	 & \keyword{SCONJ}: subordination conjunction &  \\
	
\end{tblr}

\end{frame}

\subsection{Treebanks}

\begin{frame}[fragile]
\frametitle{PoS tagging: Task description}
\framesubtitle{Treebanks}


\begin{figure}
	\begin{tcolorbox}[colback=white, colframe=blue, boxrule=1pt, text width=.9\textwidth]
		\small
\begin{alltt}
Battle-tested\keyword{/JJ} Japanese\keyword{/JJ} industrial\keyword{/JJ} managers\keyword{/NNS}
here\keyword{/RB} always\keyword{/RB} buck\keyword{/VBP} up\keyword{/RP} nervous\keyword{/JJ} newcomers\keyword{/NNS}
with\keyword{/IN} the\keyword{/DT} tale\keyword{/NN} of\keyword{/IN} the\keyword{/DT} first\keyword{/JJ} of\keyword{/IN}
their\keyword{/PP\$} countrymen\keyword{/NNS} to\keyword{/TO} visit\keyword{/VB} Mexico\keyword{/NNP} ,\keyword{/,}
a\keyword{/DT} boatload\keyword{/NN} of\keyword{/IN} samurai\keyword{/FW} warriors\keyword{/NNS} blown\keyword{/VBN}
ashore\keyword{/RB} 375\keyword{/CD} years\keyword{/NNS} ago\keyword{/RB} .\keyword{/.}
\end{alltt}
\end{tcolorbox}
\caption{Example from Penn TreeBank \cite{2003-taylor}}
\end{figure}

\begin{itemize}
	\item Universal TreeBank (\url{https://universaldependencies.org/}) \cite{2012-petrov-al}
	\begin{itemize}
		\item Open community project
		\item It aims to provide annotated corpora for several languages
		\item It uses the same PoS for all languages
	\end{itemize}
\end{itemize}

\end{frame}

\subsection{Difficulties and tools}

\begin{frame}
\frametitle{PoS tagging: Task description}
\framesubtitle{Difficulties and tools: Difficulties}

\begin{itemize}
	\item \optword{Ambiguity}
	\begin{itemize}
		\item There are words having several PoS
		\item E.g. \expword{We can can the can}
	\end{itemize}
	\item \optword{Unknown words}
	\begin{itemize}
		\item Out of vocabulary words (from the training corpus)
		\item Languages are always in evolution (creating new words and abbreviations)
	\end{itemize}
	\item \optword{Tag inconsistency}
	\begin{itemize}
		\item When there are several tags, we can find inconsistencies in the annotation
		\item E.g. \expword{In ``Brown" and ``WSJ" corpora, the word ``to" has the tag ``TO". In ``Switchboard" corpus, it has the tag ``TO" when it indicates the infinitive, otherwise ``IN" when it is a preposition.}
	\end{itemize}
\end{itemize}

\end{frame}

\begin{frame}
\frametitle{PoS tagging: Task description}
\framesubtitle{Difficulties and tools: Tools}

\begin{itemize}
	\item \url{https://nlp.stanford.edu/software/tagger.shtml}
	\item \url{https://github.com/sloria/textblob}
	\item \url{https://www.nltk.org/api/nltk.tag.html}
	\item \url{https://spacy.io/}
	\item \url{https://github.com/clips/pattern}
\end{itemize}

\end{frame}

\begin{frame}
\frametitle{PoS tagging: Task description}
\framesubtitle{Some humor}

\begin{center}
	\vgraphpage{humor/humor-identify.jpg}
\end{center}

\end{frame}

%===================================================================================
\section{Approaches}
%===================================================================================

\begin{frame}
\frametitle{Part-of-speech tagging}
\framesubtitle{Approaches}

\begin{itemize}
	\item Manual Rules
	\item Machine Learning
	\begin{itemize}
		\item Learn the rules: \expword{Transformation-Based Learning (TBL)}
		\item Statistical Models: \expword{Hidden Markov Model, Maximum Entropy Based Markov Model, Conditional Random Fields (CRF)}
		\item Neural Networks: \expword{RNNs, Transformers}
	\end{itemize}
\end{itemize}

\end{frame}

\subsection{Hidden Markov Model (HMM)}

\begin{frame}[fragile]
\frametitle{PoS tagging: Approaches}
\framesubtitle{Hidden Markov Model (HMM)}

\begin{itemize}
	\item $Q = \{t_1, t_2, \ldots, t_N\}$: tags are considered as states
	\item $A$: Transitions probability matrix where $\sum_j a_{ij} = 1,\, \forall i$
	\item $W = w_1 w_2 \ldots w_T$ : words sequence
	\item $B = b_i(w_t)$: Observation probabilities (\keyword{Emission probability}); probability of generating a word $w_t$ in a state $q_i \in Q$
	\item $\pi = [\pi_1, \pi_2, \ldots, \pi_n ]$: initial states probability distribution, where $\sum_i \pi_i = 1$
	
	\item $\hat{t} = \arg\max\limits_t P(t | w) = \arg\max\limits_t \frac{P(w|t) P(t)}{P(w)} \approx \arg\max\limits_t P(t) P(w|t)$
	\item $P(t) = P(t_1 \ldots t_T) \approx \pi_{t_1} \prod\limits_{i=1}^N P(t_i|t_{i-1}) $ where $P(t_i|t_{i-1}) \in A$ (Bi-gram) 
	\item $P(w|t) = P(w_1 \ldots w_T|t_1 \ldots t_T) = \prod\limits_{i=1}^N P(w_i|t_i) $ where $P(w_i|t_i) \in B$
\end{itemize}



\end{frame}

\begin{frame}
\frametitle{PoS tagging: Approaches}
\framesubtitle{HMM: Example}

\begin{exampleblock}{Example of a training corpus}
	\begin{itemize}
		\item a/DT computer/NN can/VB help/VB you/PN
		\item he/PN wants/VB to/TO help/VB you/PN
		\item he/PN wants/VB a/DT computer/NN
		\item he/PN can/VB swim/VB
	\end{itemize}
\end{exampleblock}

\begin{itemize}
	\item $P(VB | PN) = \frac{C(PN,\ VB)}{C(PN)} = \frac{3}{5}$, 
		  $P(VB | VB) = \frac{C(VB,\ VB)}{C(VB)} = \frac{2}{7}$,
	\item $P(he | PN) = \frac{C(PN,\ he)}{C(PN)} = \frac{3}{5}$,
	      $P(can | VB) = \frac{C(VB,\ can)}{C(VB)} = \frac{2}{7}$,
		  $P(help | VB) = \frac{C(VB,\ help)}{C(VB)} = \frac{2}{7}$,
	\item $\pi_{PN} = \frac{C(PN,\ <s>)}{C(<s>)} = \frac{3}{4} $
	\item $P(PN\ VB\ VB | he\ can\ help) \approx (\pi_{PN} P(VB | PN) P(VB | VB)) (P(he | PN) P(can | VB) P(help | VB)) $
\end{itemize}

\end{frame}

\begin{frame}
	\frametitle{PoS tagging: Approaches}
	\framesubtitle{HMM: Exercise}
	
	\begin{itemize}
		\item Use the previous corpus (example)
		\item Calculate $A$, $B$ and $Pi$ (train a HMM model)
		\item Calculate the approximate probability $P(PN\ VB\ VB | he\ can\ help)$
		\item Using greedy decoding, find tag sequence of the sentence \expword{he can help}
		\item Redo the same thing using Viterbi decoding
		\item Redo the same thing using Beam Search (K=3)
	\end{itemize}
	
\end{frame}

\subsection{Maximum entropy Markov model}

\begin{frame}
\frametitle{PoS tagging: Approaches}
\framesubtitle{Maximum entropy Markov model (MEMM)}

\begin{itemize}
	\item Let us denote $x_{n}^{m} = x_n x_{n+1} \ldots x_{m-1} x_m$
	\item Given  
	\begin{itemize}
		\item a words sequence: $w = w_1 w_2 \ldots w_n$
		\item a set of feature functions $f$ defined on a window of words $w_{i-l}^{i+l}$ and some past tags $t_{i-k}^{i-1}$
	\end{itemize}
	\item $\hat{t} = \arg\max\limits_t P(t | w) \approx \arg\max\limits_t \prod\limits_{i}  P(t_i | w_{i-l}^{i+l}, t_{i-k}^{i-1})$
	\item $P(t_i | w_{i-l}^{i+l}, t_{i-k}^{i-1}) = Softmax(\sum_j \theta_j f_j(t_i, w_{i-l}^{i+l}, t_{i-k}^{i-1}))$
\end{itemize}

\end{frame}

\begin{frame}
\frametitle{PoS tagging: Approaches}
\framesubtitle{MEMM: Features}

\begin{itemize}
	\item $w_i$ (considering words as features)
	\item previous tags 
	\item $w_i$ contains a prefix from a list ($len(prefix) \le 4$) 
	\item $w_i$ contains a suffix from a list  ($len(suffix) \le 4$) 
	\item $w_i$ contains a number 
	\item $w_i$ contains un uppercase letter
	\item $w_i$ contains a hyphen 
	\item $w_i$ is completely in uppercase
	\item $w_i$'s template (E.g. \expword{X.X.X}) 
	\item $w_i$ is uppercase with a hyphen and a number (E.g. \expword{CFC-12}) 
	\item $w_i$ is uppercase followed by words: Co., Inc., etc. after a maximum of 3 words
\end{itemize}

\end{frame}

\begin{frame}
	\frametitle{PoS tagging: Approaches}
	\framesubtitle{MEMM: Exercise}
	
	\begin{itemize}
		\item Use the previous corpus (HMM's example)
		\item Let us use the following features:
		\begin{itemize}
			\item Past, current and future words encoded using ordinal encoding (alphabetical order)
			\item Past and current PoS encoded using ordinal encoding (alphabetical order)
			\item Current word having a suffix "s" (1 or 0)
		\end{itemize}
		\item Manually train a MEMM model using these specifications:
		\begin{itemize}
			\item Parameters are initialized to 1; no bias
			\item Instead of \textbf{Softmax} use \textbf{Average normalization} as activation function
			\item Instead of \textbf{cross entropy} loss, use $J = -\sum_{i=1}^{M} t_i \hat{t}_i$
			\item Learning rate $\alpha=1$ 
			\item 3 iterations
		\end{itemize}
	\end{itemize}
	
\end{frame}

\subsection{Neural model}

\begin{frame}
\frametitle{PoS tagging: Approaches}
\framesubtitle{Neural model: Simple RNN}

\hgraphpage{rnn-simple.pdf}

\end{frame}

\begin{frame}
\frametitle{PoS tagging: Approaches}
\framesubtitle{Neural model: Stacked RNN}

\hgraphpage{rnn-stack.pdf}

\end{frame}

\begin{frame}
\frametitle{PoS tagging: Approaches}
\framesubtitle{Neural model: Bidirectional RNN}

\hgraphpage{rnn-bi.pdf}

\end{frame}

\begin{frame}
\frametitle{PoS tagging: Approaches}
\framesubtitle{Neural model: Character-based embedding RNN}

\hgraphpage{rnn-char.pdf}

\end{frame}

\begin{frame}
	\frametitle{PoS tagging: Approaches}
	\framesubtitle{Neural model: Transformers}
	
	\hgraphpage{transformers.pdf}
	
\end{frame}

\begin{frame}
\frametitle{PoS tagging: Approaches}
\framesubtitle{Neural model: Some humor}

\begin{center}
	\vgraphpage{humor/humor-learn.jpg}
\end{center}

\end{frame}

\insertbibliography{NLP05}{*}

\end{document}

